% =============================================================================
% SECTION 2: CONTACT, NEGATION HIERARCHY, AND PARTITION DEPTH MINIMISATION
% =============================================================================

\section{Contact as Partition Depth Minimisation}
\label{sec:contact}

\subsection{The Hierarchy of Negation}
\label{subsec:neg_hierarchy}

Postulate~\ref{post:negation} establishes that items exist by negation against
the totality $\mathcal{U}$. This is the \emph{global negation} form of
existence: $\mathcal{I}$ is defined as not-$\mathcal{U}$. Global negation is
structurally expensive: the boundary separating $\mathcal{I}$ from all of
reality must be maintained against an unbounded complement. The boundary
thickness $\beta$ must span the full partition depth required to distinguish
$\mathcal{I}$ from every other item simultaneously.

A strictly weaker but physically realised form of negation arises when two
items $\mathcal{I}_1$ and $\mathcal{I}_2$ are in proximity: each begins to be
defined not merely against $\mathcal{U}$ but against the other. We call this
\emph{local negation}.

\begin{definition}[Local Negation]
\label{def:local_neg}
Items $\mathcal{I}_1$ and $\mathcal{I}_2$ are in \emph{local negation} when
the boundary between them has smaller thickness than the boundary between either
item and the global complement $\overline{\mathcal{I}_1 \cup \mathcal{I}_2}$.
Formally:
\begin{equation}
    \beta(\mathcal{I}_1, \mathcal{I}_2) < \beta(\mathcal{I}_1, \mathcal{U})
    \quad \text{and} \quad
    \beta(\mathcal{I}_2, \mathcal{I}_2) < \beta(\mathcal{I}_2, \mathcal{U})
    \label{eq:local_neg}
\end{equation}
\end{definition}

Local negation is more efficient than global negation in the following precise
sense: it requires a smaller residue $\beta$ to maintain an unambiguous
distinction. An item in local negation with a specific partner exists more
precisely than an item in only global negation against all of reality, because
the local boundary is maintained against a single finite complement rather than
an infinite one.

\subsection{Contact as Minimum Viable Partition}
\label{subsec:contact_min}

\begin{definition}[Contact]
\label{def:contact}
Two items $\mathcal{I}_1$ and $\mathcal{I}_2$ are in \emph{contact} when their
mutual boundary thickness equals the floor:
\begin{equation}
    \beta(\mathcal{I}_1, \mathcal{I}_2) = \beta_{\min}
    \label{eq:contact}
\end{equation}
Contact is the minimum viable partition: the state in which the two items are
maximally locally negated while remaining distinct.
\end{definition}

\begin{theorem}[Contact as Existence Optimum]
\label{thm:contact_optimum}
For two items $\mathcal{I}_1$ and $\mathcal{I}_2$, the contact state minimises
the total boundary thickness required to maintain both items as distinct entities.
Any state with $\beta(\mathcal{I}_1, \mathcal{I}_2) > \beta_{\min}$ has greater
total boundary thickness and is therefore a less efficient state of mutual
existence.
\end{theorem}

\begin{proof}
The total boundary cost of maintaining $\mathcal{I}_1$ and $\mathcal{I}_2$ as
distinct entities consists of:
\begin{enumerate}
    \item The mutual boundary $\beta(\mathcal{I}_1, \mathcal{I}_2)$
    \item The global boundary of their union against $\overline{\mathcal{I}_1
          \cup \mathcal{I}_2}$
\end{enumerate}
When $\beta(\mathcal{I}_1, \mathcal{I}_2) > \beta_{\min}$, there exist
intervening partition layers between the two items. Each layer is an item in its
own right, requiring its own global boundary against $\mathcal{U}$. The total
boundary cost therefore includes the costs of all intervening layers. As
$\beta(\mathcal{I}_1, \mathcal{I}_2)$ decreases toward $\beta_{\min}$,
intervening layers are eliminated and the total boundary cost decreases
monotonically. At $\beta_{\min}$, no intervening layers remain and the boundary
cost is minimised.
\end{proof}

\subsection{Partition Depth Between Items}
\label{subsec:depth_between}

\begin{definition}[Inter-Item Partition Depth]
\label{def:inter_depth}
The \emph{partition depth} $D(\mathcal{I}_1, \mathcal{I}_2)$ between two items
is the number of distinct partition boundaries that must be traversed to reach
$\mathcal{I}_2$ from $\mathcal{I}_1$ through the contact chain. Equivalently,
it is the minimum number of intervening items in the shortest contact chain
connecting $\mathcal{I}_1$ to $\mathcal{I}_2$.
\end{definition}

At contact, $D(\mathcal{I}_1, \mathcal{I}_2) = 1$. In general, $D$ satisfies
the triangle inequality:
\begin{equation}
    D(\mathcal{I}_1, \mathcal{I}_3) \leq D(\mathcal{I}_1, \mathcal{I}_2) +
    D(\mathcal{I}_2, \mathcal{I}_3)
    \label{eq:triangle}
\end{equation}
making partition depth a metric on the set of items.

\subsection{The Partition Depth Minimisation Principle}
\label{subsec:pdm}

\begin{postulate}[Partition Depth Minimisation]
\label{post:pdm}
In the absence of constraints preventing contact, items in bounded phase space
evolve toward states of minimum inter-item partition depth. This is not a
force law but a consequence of the efficiency of local negation: items move
toward contact because contact is the most efficient state of mutual existence.
\end{postulate}

This postulate replaces all force laws. We now derive each of the four
fundamental interactions as a consequence of partition depth minimisation
operating at different scales and different partition depths.

\subsection{Gravity as Partition Depth Reduction}
\label{subsec:gravity}

Consider two macroscopic items $\mathcal{I}_1$ and $\mathcal{I}_2$ separated by
a medium $\mathcal{M}$ of $n$ partition layers. The total partition depth is
$D(\mathcal{I}_1, \mathcal{I}_2) = n + 2$ (the $n$ layers of $\mathcal{M}$
plus the two boundary layers at each surface).

The evolution of $D$ over time follows from the monotone increase of $M$. As
$M$ increases, the system explores available partition configurations. The
configurations with lower $D$ are preferred because they require less total
boundary maintenance. The rate of change of $D$ is:
\begin{equation}
    \frac{dD}{dt} = -\frac{\partial \beta_{\mathrm{total}}}{\partial D} \cdot
    \frac{dM/dt}{|\mathcal{P}(D)|}
    \label{eq:depth_reduction}
\end{equation}
where $|\mathcal{P}(D)|$ is the number of partition states at depth $D$ and
$\beta_{\mathrm{total}}$ is the total boundary thickness of the system.

The observable consequence of Eq.~\eqref{eq:depth_reduction} is what we call
gravitational attraction. Two massive items --- items with large internal
partition count $M$ --- have large $|\mathcal{P}(D)|$ and therefore large
$dD/dt$. They reduce partition depth more rapidly. This is the origin of the
proportionality of gravitational force to mass: mass is internal partition count
(as shown in Section~\ref{sec:emergence}), and larger partition count generates
faster depth reduction.

\begin{theorem}[Gravity from Partition Depth]
\label{thm:gravity}
The acceleration of $\mathcal{I}_1$ toward $\mathcal{I}_2$ is:
\begin{equation}
    a = -G \frac{M_2}{r^2}
    \label{eq:gravity}
\end{equation}
where $G$ is the \emph{contact geometry constant} --- the conversion factor
between partition depth gradient and spatial acceleration --- $M_2$ is the
partition count of $\mathcal{I}_2$, and $r$ is the spatial distance
(itself a consequence of partition depth, as derived in
Section~\ref{sec:emergence}).
\end{theorem}

\begin{remark}[The Irreducibility of $G$]
\label{rem:G_irreducible}
The constant $G$ is irreducible to partition counts. Every other fundamental
constant ($c$, $\hbar$, $k_B$) is expressible as a ratio of partition counts
and tick energies, because they describe rates of partition processes. $G$
describes the relationship between the partition depth gradient and the spatial
geometry that partition depth induces. Spatial geometry is the contact ground
state --- the configuration the universe reaches when partition depth is
minimised. The ground state is prior to the process that approaches it; therefore
$G$ cannot be expressed in terms of the process. This is why $G$ is the only
fundamental constant that resists reduction in the partition framework, and why
its measurement is uniquely difficult: every other constant can be pinned by
counting; $G$ can only be measured by comparing contact geometries directly.
\end{remark}

\subsection{The Contact Chain and Structural Self-Reinforcement}
\label{subsec:chain}

When item $\mathcal{I}_1$ comes into contact with $\mathcal{I}_2$, and
$\mathcal{I}_2$ is already in contact with $\mathcal{I}_3$, the following
occurs. The contact $\mathcal{I}_1$--$\mathcal{I}_2$ establishes a new local
negation relationship: $\mathcal{I}_1$ is now *not* $\mathcal{I}_2$
unambiguously. But this new relationship also sharpens the existing boundary
between $\mathcal{I}_2$ and $\mathcal{I}_3$: $\mathcal{I}_2$ is now not only
*not* $\mathcal{I}_3$ but *not* $\mathcal{I}_3$-while-in-contact-with-$\mathcal{I}_1$.
The contact chain acquires additional specificity at every node.

\begin{theorem}[Chain Self-Reinforcement]
\label{thm:chain_reinforce}
For a contact chain $\mathcal{I}_1 - \mathcal{I}_2 - \cdots - \mathcal{I}_n$,
adding a new contact $\mathcal{I}_0 - \mathcal{I}_1$ increases the definitional
precision of every item in the chain. Specifically, the partition depth of the
chain from $\mathcal{I}_0$ to $\mathcal{I}_n$ satisfies:
\begin{equation}
    D(\mathcal{I}_0, \mathcal{I}_n) = n
    \label{eq:chain_depth}
\end{equation}
and the addition of $\mathcal{I}_0$ adds one unit of definitional specificity
to each of $\mathcal{I}_1, \ldots, \mathcal{I}_n$ without adding any units of
global negation cost to any of them.
\end{theorem}

\begin{proof}
Before $\mathcal{I}_0$ is added, the chain $\mathcal{I}_1 - \cdots -
\mathcal{I}_n$ has each item defined by its local negation partners within the
chain plus its global negation against $\mathcal{U}$. After $\mathcal{I}_0$
contacts $\mathcal{I}_1$, the item $\mathcal{I}_1$ gains a new local negation
partner. The boundary $\beta(\mathcal{I}_1, \mathcal{I}_0) = \beta_{\min}$ is
established without changing any other boundary in the chain, since the
boundaries $\beta(\mathcal{I}_k, \mathcal{I}_{k+1})$ are already at $\beta_{\min}$
(by definition of contact). The only change is that the \emph{context} of each
item in the chain now includes one additional specific partner at $\mathcal{I}_1$.
This context propagates through the chain as causal residues, increasing
definitional precision at each node.
\end{proof}

\subsection{Normal Force as Floor Resistance}
\label{subsec:normal}

The contact state $\beta(\mathcal{I}_1, \mathcal{I}_2) = \beta_{\min}$ is
stable because any further reduction would require $\beta < \beta_{\min}$, which
is forbidden by Postulate~\ref{post:floor}. Attempting to drive the two items
below $\beta_{\min}$ requires introducing a restoring tendency --- the resistance
of the partition floor to being violated.

This restoring tendency is what classical mechanics calls the normal force. It
is not a separate force. It is the gradient of the floor:
\begin{equation}
    F_{\mathrm{normal}} = -\frac{\partial \beta_{\min}}{\partial D}
    \bigg|_{D \to 0^+}
    \label{eq:normal}
\end{equation}
The normal force is zero when $D > D_{\min}$ (the floor is not being approached)
and becomes unbounded as $D \to 0^+$ (the floor is impenetrable). This explains
why the normal force appears only at contact and not before: it is not a
pre-existing repulsion but the gradient of the floor, which is only encountered
when partition depth reaches its minimum.

\begin{remark}[Craters and Deformation]
\label{rem:craters}
When a stone strikes the ground at high velocity, the impact carries accumulated
partition count $\Delta M = \dot{M} \cdot v \cdot t$ in the direction of motion.
This partition count must be redistributed across the contact chain when the
floor is reached. The redistribution manifests as deformation: the ground's
partition structure is rearranged to accommodate the incoming partition count.
The crater is the residue of this rearrangement --- the new partition depth
minimum of the combined stone-ground system after the incoming $\Delta M$ has
been absorbed. The fact that craters form \emph{after} impact, rather than the
ground pre-emptively resisting, follows immediately: the floor is encountered at
the moment of contact, not before.
\end{remark}
